\section{Certified square-root frontiers and a subexponential scalar filter}
\label{sec:separator-orders}

The previous integration reduced the cost of observing a fixed scan.  We now
address its order.  The maintained Python implementation constructs a crossing
order of width $O(\sqrt n)$ for every closed classical projection with $n$
crossings.  It supplies a separately checkable hierarchy and compares the
result with the existing greedy order, retaining the smaller measured frontier.
Combined with the exact fixed-color Potts evaluator, this gives a general
$\operatorname{poly}(n)2^{O(\sqrt n)}$ bit bound for \emph{one scalar Jones
specialization}, with optional resource caps disabled.  It does not give that
bound for unknot recognition: equality with the unknot's scalar is inconclusive.

This is an implementation of classical planar-separator ideas
\cite{lipton-tarjan}, not a new planar-separator theorem.  The contribution here
is a constructive order, a graph-only certificate verifier, its integration
with the existing exact scalar backend, and explicit accounting of which
recognition bottleneck remains.  Nothing in this section requires connected
prefixes, a common boundary disk, or a bound on surviving Khovanov objects.

\subsection{A bounded-radius middle band}

Represent the projection by a plane multigraph $G$: crossings are vertices,
diagram arcs are edges, and the PD rows specify the cyclic orders of half-edges.
The degree is four, with loops counted twice.  The constructor checks that each
edge label occurs twice and that the rotation system is spherical on every
connected component.  Loops may then be deleted for separator construction:
they neither connect distinct components nor cross a cut of crossing vertices.
Parallel edges remain in the embedded map.  Disconnected projections are
handled component by component.

Consider one connected recursive subproblem with $s$ vertices.  Take breadth
first search levels $V_0,\ldots,V_d$ from a root, and put
$r=\lceil\sqrt{s}\rceil$.  Let $a$ be the first level for which
$\sum_{i\le a}|V_i|\ge s/2$.  If $a\ge r$, choose a smallest level in
$[a-r,a-1]$ as the lower boundary $\ell$; otherwise use the empty virtual
level $\ell=-1$.  If $a+r\le d$, choose a smallest level in
$[a+1,a+r]$ as the upper boundary $u$; otherwise put $u=d+1$.
Each real chosen level has size at most $s/r\le r$.

The vertices strictly below $\ell$ and strictly above $u$ number at most
$s/2$ on either side.  An edge changes its BFS distance by at most one, so
removing the two boundary levels separates these outside vertices from the
middle band.  Retain levels below $u$, and, when $\ell\ge0$, contract the
connected inner BFS tree on levels at most $\ell$ to the root.  Delete the
resulting loops.  The resulting plane multigraph $H$ has radius at most
$2r$.  Give weight one to each original middle-band vertex and weight zero
to the contracted root.  The total weight is at most $s$.

The ribbon-map implementation contracts a tree by following its boundary
through the removed half-edges.  It does not repeatedly rebuild the graph
after individual edge contractions.  This matters for long bands and large
inner balls.  An independent check after construction verifies the actual
separator size and every remaining component's size in the original graph.

\subsection{Three short paths from a dual centroid}

For each face of $H$, add a center and a spoke to every corner occurrence.
Keep the old edges.  A bridge or a repeated vertex on a face walk therefore
causes repeated spokes, not an assumption that the face boundary is a simple
cycle.  The resulting map is triangulated; every new center has distance at
most $2r+1$ from the root.  The one-vertex case is handled directly.

Choose a BFS spanning tree $T$ of this triangulation.  The dual edges of
primal edges outside $T$ form a dual spanning tree $T^*$.  Assign each weighted
original vertex to one incident triangular face.  Choose a weighted centroid
face $f$ of $T^*$: every component of $T^*-f$ has at most half the total
assigned weight.  Remove the union $P$ of the three primal tree paths from
the root to the vertices of $f$.

To see balance, consider a dual edge incident to $f$.  Its primal edge and
the tree path between its endpoints bound a fundamental cycle.  The cycle's
vertices lie in $P$, since both endpoints belong to $f$.  It separates the
corresponding branch of the dual cotree from the other branches.  Consequently
each surviving primal component has its assigned faces in one branch of
$T^*-f$.  Any weight assigned to $f$ itself belongs to a vertex of $f$, hence
to $P$.  Every remaining component therefore has at most half the middle-band
weight.  This argument permits arbitrary nonnegative face weights; it does
not require each individual assigned weight to be small.

Return only the \emph{original middle-band vertices} of $P$, together with
the two boundary levels.  Discarding artificial centers from this separator
is safe: paths in the original graph use old vertices and old edges only.
The old edges were retained during triangulation, so any original path
avoiding the returned vertices is also a path in the triangulation avoiding
$P$.  The contracted root represents vertices already separated from the
middle by the lower level.  Thus the original graph has a separator $S$ with
\[
 |S|\le 2r+3(2r+2)=8\lceil\sqrt{s}\rceil+6,
 \qquad |C|\le s/2\quad(C\text{ a component of }G-S).
\]
The constants are deliberately loose.  Small separators on practical inputs
matter more than optimizing this universal numerical bound.

\subsection{From the hierarchy to a certified order}

Recursively order each component of $G-S$, then put $S$ last.  Stop at a
fixed leaf size, eight by default; the API permits integer cutoffs from one
through 64.  Process distinct initial connected components consecutively.
At a prefix cut, any crossing arc is incident to a separator on the current
recursion path, or to a vertex of the current leaf.  Completed sibling
components have no edges to unvisited siblings.  Since the original degree
is four, the frontier is at most
\[
 4\left(\sum_{S\text{ on the path}}|S|+|\text{leaf}|\right).
\]
Sizes halve along a path.  The square-root terms form a convergent geometric
sum, while the additive constants contribute $O(\log n)$, giving width
$O(\sqrt n)$.  This holds for every prefix, including prefixes inside a
separator's own vertex list.

The JSON certificate records the input digest, leaf cutoff, separator and
child lists, flattened order, actual maximum and total frontier, and the
recursively computed bound
\[
 b(\text{leaf})=4|\text{leaf}|,\qquad
 b(\text{node})=4|S|+\max_C b(C).
\]
Each crossing occurs once in the hierarchy.  With normalized crossing
indices the hierarchy uses $O(n\log n)$ bits; arbitrary input-label bit
length also counts toward input processing.  Construction and verification
take polynomial time.  The verifier rebuilds adjacency directly from edge
owners and checks coverage, disjointness, the exact remaining components,
half-balance, separator sizes, order, actual profile and the bound.  It calls
no planar-map or separator constructor.  Thus even a mistake in the embedded
construction cannot silently publish an incorrect certified frontier bound.

The practical wrapper also computes or accepts a candidate order.  It
chooses the lexicographically smaller pair (maximum frontier, total frontier)
between that order and the separator order.  If greedy wins, the separator
witness is retained: explicit profile comparison proves that the selected
order has the same asymptotic upper bound.  This is a bounded-choice ordering
policy, not a progress-based restart of a scan.  It complements the earlier
adaptive elimination and observer fallbacks.

\subsection{Exactly which subexponential computation follows}

For a fixed checkerboard shading, let $f_i$ count Tait vertices incident to
both processed and unprocessed crossings, and let $w_i$ be the crossing
frontier.  The corner argument in Section~\ref{sec:potts-integration} gives
$2f_i\le w_i$, including repeated face walks.  A fixed $q\ge5$ Potts query
uses at most $q^{f_i+2}$ temporary assignments and exact pairs in
\[
 R_q=\mathbb Z[x]/(x^2-(q-2)x+1).
\]
Coefficient bit lengths are $O(n\log q)$; arithmetic and boundary-key work
add polynomial factors.  Substituting the certified width yields
\[
 T_{\mathrm{scalar}}(n)=\operatorname{poly}(n)\,2^{O(\sqrt n)}
 \quad\text{for fixed }q.
\]
This is a bit-complexity statement, not a count of unit-cost operations on
unbounded integers.  For the existing modular matching filter, allowing all
boundary pairings gives the weaker but still subexponential bound
$\operatorname{poly}(n)2^{O(\sqrt n\log(n+1))}$ on the same order.
These are universal upper bounds for scalar evaluation, not claims that
our implementation or constants are optimal.

The optional \code{potts-separator} Jones backend constructs the witness and
then calls the existing exact evaluator.  With the Python API, setting both
\code{max\_states=None} and \code{max\_transitions=None} disables the local
caps for the theoretical complete query.  Production defaults retain the
4096-state and 200\,000-transition caps.  Exhaustion is inconclusive, and
the recognition deadline remains global.  A zero local budget declines
before constructing a hierarchy.  Once a complete certificate has been
published, scalar exhaustion retains its order for the homology fallback.
If subsequent Reidemeister reductions change crossing indices, the witness
is rebuilt on the reduced diagram before supplying an order to a window
query or the scanner.  A global interruption during preparation publishes
no partial certificate.

Three obstacles prevent promoting this to a general recognition theorem.
First, agreement at a scalar specialization is not an unknot certificate.
Second, a frontier bound alone does not bound the number of surviving
relative Khovanov objects or intermediate matrix size.  Third, this hierarchy
does not certify common-disk prefixes, so the stronger conditional bounds of
the disk-transfer and nice-window backends cannot simply be substituted.
The descending grid family already described in this report has raw width
$\Omega(\sqrt n)$ for every order.  That obstructs a universal
polylogarithmic-width strategy on a fixed diagram, but it is no lower bound
for recognition: a descending certificate immediately handles the family.

\subsection{Validation and measurements}

All 589 integrated tests pass.  The new tests compare acceptance of 500 random
ribbon maps with an independent spherical-map implementation, verify
certificate corruption is rejected, and exercise loops, disconnected
projections, long bands and grids.  Sixteen small knot diagrams are checked
against independent cube homology using actual separator orders; both Potts
shadings agree with the previous exact evaluator.  Resource tests distinguish
partial preparation from a completed certificate followed by scalar
exhaustion, and check that a changed diagram receives a fresh certificate.

The separate reproducible driver \path{fast/audit_separator_orders.py}
audits 114 inputs, including 100 randomly generated, crossing-permuted and
edge-relabelled braid closures.  Its independent hierarchy checker uses
union-find over edge-owner pairs rather than the production component
traversal.  It independently counts cut edges and checks balance, coverage,
components and the stored bound.  The result archive retains all inputs,
certificates, source hashes, five shuffled ordering rounds and three shuffled
kernel rounds, with warm-ups excluded.  Table~\ref{tab:separator-orders}
separates measured frontiers from preparation costs.

\begin{table}[ht]
\centering\small
\begin{tabular}{@{}lrrrrrr@{}}
\toprule
Family & $n$ & Greedy & Separator & Selected & Greedy ms & Bounded ms \\
\midrule
Grid $8\times8$ & 64 & 10 & 46 & 10 & 0.626 & 4.524 \\
Grid $32\times32$ & 1024 & 34 & 214 & 34 & 6.082 & 142.782 \\
Tree medial, height 5 & 62 & 16 & 12 & 12 & 1.118 & 4.492 \\
Tree medial, height 7 & 254 & 64 & 14 & 14 & 5.218 & 26.624 \\
Tree medial, height 9 & 1022 & 256 & 18 & 18 & 23.653 & 140.075 \\
\bottomrule
\end{tabular}

\caption{Widths and median order-preparation times. Bounded preparation
includes the separator witness, its verification and the greedy comparison;
these are additional costs, not speedups.}\label{tab:separator-orders}
\end{table}

The tree-medial family consists of medial diagrams of embedded complete
binary trees, with one crossing per tree edge.  These are simple nugatory
unknot diagrams used to stress the greedy order; they are not hard
recognition instances.  In isolated modular matching queries, height five
falls from a median 47.665\,ms to 1.941\,ms (24.6 times), excluding order
preparation.  At heights six and seven, all greedy runs exhaust the state
cap while the selected orders finish in 4.126 and 8.788\,ms.  Those censored
runs have no completion-time speedup ratio.  All completed scalar answers
are inconclusive for recognition.

The exact Potts queries on these same trees differ little between orders:
the advantageous checkerboard shading already has trivial frontier.  At
height seven the medians are 2.572 and 2.585\,ms, far below the extra order
preparation cost.  The grids also retain their better greedy order.  These
results support keeping the new option explicit: the certified asymptotic
bound and a reduction in matching-filter work do not establish a general
practical speedup, a full-recognition speedup, or a smaller homology complex.
